% Part:sets-functions-relations
% Chapter: size-of-sets
% Section: equinumerous-sets

\documentclass[../../../include/open-logic-section]{subfiles}

\begin{document}

\olfileid{sfr}{siz}{equ}

\olsection{तुल्यबलता}

संचांच्या ``आकारमाना''ची एक अंतःप्रज्ञ कल्पना आपल्याकडे आहे आणि ती
सांत संचांसाठी व्यवस्थित चालते. पण अनंत संचांचे काय? \emph{कोणत्याही}
आकारमानाच्या दोन संचांच्या आकारमानांची औपचारिक रीतीने तुलना करायची
असेल, तर दोन संच एकाच आकारमानाचे कधी असतात याची व्याख्या प्रथम करणे
योग्य ठरते. फ्रेगे यांचे पुढील उदाहरण पाहा:
\begin{quote}
  उपाहारगृहातील कर्मचाऱ्याला मेजावर ताटांइतक्याच सुऱ्या मांडल्या आहेत
  याची खात्री करायची असेल, तर प्रत्येक ताटाच्या उजवीकडे एक सुरी ठेवून
  आणि मेजावरील प्रत्येक सुरी कोणत्या तरी ताटाच्या उजवीकडे येईल असे
  केल्यास त्याला ताटे किंवा सुऱ्या मोजण्याची गरज नसते. अशा प्रकारे
  ताटे आणि सुऱ्या एकास-एक संगतीने परस्परांशी जोडली जातात, आणि तीही
  त्याच अवकाशीय संबंधाने. \citep[\S70]{Frege1884}
\end{quote}
या उताऱ्यातील अंतर्दृष्टी पुढील औपचारिक व्याख्येत मांडता येते:

\begin{defn}\ollabel{comparisondef}
  $A$ हा $B$ शी \emph{तुल्यबल} असतो, आणि ते $\cardeq{A}{B}$ असे लिहितो,
  तेव्हा आणि केवळ तेव्हाच, जेव्हा $f \colon A \to B$ असे
  !!a{bijection} असते.
\end{defn}

\begin{prop}\ollabel{equinumerosityisequi}
तुल्यबलता हा तुल्यता संबंध आहे.
\end{prop}

\begin{proof} 
तुल्यबलता परावर्ती, सममित आणि संक्रमक आहे, हे दाखवावे लागेल. $A, B$
आणि $C$ हे संच असू द्या.

\emph{परावर्तिता.} प्रत्येक $x \in A$ साठी $\Id{A} (x) = x$ असे
असलेले एकरूपता फलन $\Id{A} \colon A \to A$ हे !!a{bijection} आहे.
म्हणून $\cardeq{A}{A}$.

\emph{सममितता.} $\cardeq{A}{B}$ आहे, म्हणजे $f\colon A \to B$ असे
!!a{bijection} आहे, असे समजा. $f$ हे !!{bijective} असल्याने त्याचे
व्युत्क्रम $f^{-1}$ अस्तित्वात आहे आणि तेही !!{bijective} आहे. म्हणून
$f^{-1}\colon B \to A$ हे !!a{bijection} आहे; त्यामुळे
$\cardeq{B}{A}$.

\emph{संक्रमकता.} $\cardeq{A}{B}$ आणि $\cardeq{B}{C}$ आहे, म्हणजे
$f\colon A \to B$ आणि $g\colon B \to
C$ अशी !!{bijection}s आहेत, असे समजा. मग संयोजन
$\comp{f}{g}\colon A \to C$ हे !!{bijective} आहे; म्हणून
$\cardeq{A}{C}$.
\end{proof}

\begin{prop}
$\cardeq{A}{B}$ असेल, तर $A$ हा !!{enumerable} असतो तेव्हा आणि केवळ
तेव्हाच $B$ तसा असतो.
\end{prop}

\begin{editorial}
पुढील सिद्धतेत \olref[enm]{sec} समाविष्ट असेल, तर
\olref[enm]{defn:enumerable} वापरली जाते; अन्यथा
\olref[enm-alt]{defn:enumerable} वापरली जाते.
\end{editorial}

\begin{proof}
$\cardeq{A}{B}$ आहे, म्हणून $f \colon A
\to B$ असे काही !!{bijection} आहे, आणि $A$ हा !!{enumerable} आहे,
असे समजा.
\oliflabeldef{sfr:siz:enm:defn:enumerable}{
  मग एकतर $A = \emptyset$ असतो, किंवा $g\colon \PosInt \to A$ असे
  !!a{surjective} फलन असते. $A = \emptyset$ असेल, तर $B = \emptyset$
  सुद्धा असतो (अन्यथा एखादा !!a{element}~$y \in B$ असेल,
  पण $x \in A$ असा कोणताही घटक नसल्याने $f(x) = y$ अशी पूर्वप्रतिमा
  मिळणार नाही). दुसरीकडे, $g\colon \PosInt \to A$ हे !!{surjective}
  असेल, तर $\comp{g}{f} \colon \PosInt \to B$ हे !!{surjective} आहे.
  हे पाहण्यासाठी $y \in B$ असू द्या. $f$ हे !!{surjective} असल्याने
  $f(x) = y$ असा एखादा $x \in A$ घटक असतो. $g$ हे !!{surjective}
  असल्याने $g(n) =
  x$ अशी एखादी $n \in \PosInt$ संख्या असते. त्यामुळे
\[
(\comp{g}{f})(n) = f(g(n)) = f(x) = y
\]
  आणि म्हणून $\comp{g}{f}$ हे !!{surjective} आहे. $\comp{g}{f}$ हे~$B$
  चे प्रगणन आहे; त्यामुळे $B$ हा !!{enumerable} आहे.}
{
  मग एकतर $A  = \emptyset$ असतो, किंवा ज्याची व्याप्ती $A$ आहे आणि
  ज्याचा प्रांत $\Nat$ किंवा नैसर्गिक संख्यांचा आरंभीचा खंड आहे असे
  !!a{bijection}~$g$ असते. $A = \emptyset$ असेल, तर $B =
  \emptyset$ सुद्धा असतो (अन्यथा एखादा~$y \in B$ असेल, पण $x
  \in A$ असा कोणताही घटक नसल्याने $f(x) = y$ अशी पूर्वप्रतिमा मिळणार
  नाही). म्हणून आपल्याकडे !!{bijection}~$g$ आहे असे समजा. मग
  $\comp{g}{f}$ हे !!a{bijection} आहे; त्याची व्याप्ती~$B$ आहे आणि
  त्याचा प्रांत~$g$ च्या प्रांतासारखाच आहे (म्हणजे $\Nat$ किंवा त्याचा
  आरंभीचा खंड). म्हणून $B$ हा !!{enumerable} आहे.}
%READERNOTE{OLSIZ-006: गोठवलेल्या इंग्रजी सिद्धतेच्या दोन्ही पर्यायी शाखांतील रिक्त-संच प्रकरणात g(x)=y असे लिहिले आहे. त्या शाखेत g अजून अस्तित्वात नसते; A पासून B पर्यंतचे आधी दिलेले द्विएक फलन f वापरूनच B रिक्त असल्याचा निष्कर्ष मिळतो. मराठी मजकुरात दोन्ही ठिकाणी f(x)=y केले आहे.}
%READERNOTE{MRSIZ-004: गोठवलेल्या पर्यायी शाखेत एकाच प्रांताला आधी “initial sequence of natural numbers” आणि नंतर “initial segment” म्हटले आहे. आधीच्या औपचारिक व्याख्येशी सुसंगत म्हणून मराठीत दोन्हीकडे “आरंभीचा खंड” वापरला आहे.}

$B$ हा !!{enumerable} असेल, तर~$f$ ऐवजी
!!{bijection} $f^{-1}\colon B \to A$ घेऊन तोच युक्तिवाद पुन्हा केल्यास
$A$ हा !!{enumerable} आहे, हे मिळते.
\end{proof}

\begin{prob}
$\cardeq{A}{C}$ आणि $\cardeq{B}{D}$ असेल, आणि $A \cap B =
C \cap D = \emptyset$ असेल, तर $\cardeq{A \cup B}{C \cup D}$ आहे,
हे दाखवा.
\end{prob}

\begin{prob}
$A$ हा अनंत आणि !!{enumerable} असेल, तर $\cardeq{A}{\Nat}$ आहे, हे
दाखवा.
\end{prob}

\end{document}
